%============================================================ 5. 실험 설계
\section{실험 설계}
\label{sec:experiment}

실험의 목적은 두 계층의 기여를 분리하여 계층별로 식별하는 것이다. 이를 위해 배정 축과
기동 축을 각각 이분한 $2\times2$ 조건 행렬을 두고, 모든 조건을 같은 시드 집합으로 수집해
시드 대응 차이로 판정한다. 이 절은 그 재현에 필요한 환경·학습 실행·조건·판정 기준을 한곳에
모은다.

\subsection{실험 환경}
\label{sec:env}

시뮬레이터·관측 래스터화·평가 하네스는 자체 구현이며 공개한다(Data availability 참조).
교전 설정값은 Table~\ref{tab:scenario}, 지표는 Table~\ref{tab:metrics}를 따르고, 두 계층
모두 25\,step 주기로 결정하며 지휘관 재계획도 배치 설정과 같이 매 결정마다 수행한다. 학습·평가·
언어모델 서빙은 NVIDIA RTX 4070(12\,GB) 단일 GPU 한 대에서 이루어졌고, 지연 측정 때는
시뮬레이터·정책 추론·언어모델이 같은 GPU를 공유하였다(\S\ref{sec:res_latency}). 로컬
지휘관(Qwen2.5 7B·14B)은 Ollama로 서빙하며 14B는 Q4\_K\_M 양자화·문맥 16{,}384 토큰이고,
클라우드 지휘관(Gemini 3.5 Flash-Lite)은 상용 API를 쓴다. 세 모델 모두 온도 0, 같은
프롬프트, 같은 JSON 스키마의 구조적 출력(\S\ref{sec:schema})을 사용하며, 스키마 검증에
실패한 응답은 휴리스틱 배정으로 대체한다(폴백률은 세 지휘관 모두 0.007 이하).

\subsection{학습 실행}
\label{sec:protocol}

학습은 실해역 6곳(Fig.~\ref{fig:sites}; 육지 비율 0--7.5\,\%, 요격 구역 잠식 0--6\,\%,
최근접 육지 여유 1.6\,km--$\infty$)을 에피소드마다 돌려가며 진행하여 정책이 한 해안선의
형상을 외우지 않게 했고, 평가는 그중 통영 매물도 서측 한 곳에서
한다. 총 800 업데이트($G=8$, 후보 평가 구간 180\,step)를 돌리며 200 업데이트마다 greedy
평가를 두고, \textbf{그 최고점(600 업데이트, 포획률 0.805; 휴리스틱 기준선 0.744)의
가중치}를 채택하였다 --- 학습 곡선은 매 업데이트 새로 표집한 전장의 난이도로 진동하므로
채택 판정은 고정 프로토콜의 greedy 평가로 한다.

\subsection{조건과 수집}
\label{sec:paired}

\begin{table*}[tp]
\centering
\caption{$2\times2$ 조건 행렬. 조건 사이에 달라지는 것은 두 스위치뿐이며, 체크포인트·환경
         설정·난수 시드는 모두 동일하다.}
\label{tab:conditions}
\footnotesize
\begin{tabular}{llp{0.42\linewidth}}
\toprule
조건 & 배정 $\times$ 기동 & 무엇을 측정하는가 \\
\midrule
\texttt{heur\_heur} & 휴리스틱 $\times$ 휴리스틱 & 기준선(baseline) \\
\texttt{heur\_unet} & 휴리스틱 $\times$ U-Net    & \textbf{기동 계층의 순수 기여}
                                                   --- LLM을 쓰지 않는다 \\
\texttt{llm\_heur}  & LLM $\times$ 휴리스틱      & \textbf{전략 계층의 순수 기여} \\
\texttt{llm\_unet}  & LLM $\times$ U-Net         & 제안 체계(상호작용 포함) \\
\bottomrule
\end{tabular}
\end{table*}

휴리스틱 배정과 기동은 \S\ref{sec:heuristic}의 기준 방법이며 학습 정책과 같은 행동공간을
거치므로 네 조건은 직접 비교 가능하다(Table~\ref{tab:conditions}). 전략 계층에는 능력
수준이 다른 세 모델 --- Qwen2.5 7B, Qwen2.5 14B(로컬), Gemini 3.5 Flash-Lite(클라우드) ---
을 두어 총 8조건이 된다. 모든 조건은 동일한 시드 집합(0--29)과 세 포메이션을 사용하며,
8조건 $\times$ 3포메이션 $\times$ 30시드 $=$ 720 에피소드를 수집하였다.

\subsection{판정 기준과 계측}
\label{sec:stats}

조건 $X$ 와 기준선 $Y$ 의 시드 $k$ 지표값을 $x_k$, $y_k$ 라 할 때 대응 차이와 효과크기는
\begin{equation}
d_k = x_k - y_k, \qquad
\bar{d} = \frac{1}{n_{\mathrm{s}}}\sum_{k} d_k, \qquad
\dz = \frac{\bar{d}}{s_d}
\label{eq:paired}
\end{equation}
이며($n_{\mathrm{s}} = 30$, $s_d$ 는 $d_k$ 의 표준편차), $\bar{d}$ 의 95\,\% 신뢰구간은
BCa 부트스트랩~\citep{efron1993bootstrap, diciccio1996bca}($B = 9999$)으로 구한다. 구간이
0을 포함하면 기준선과 다르다고 말하지 않는다. 지표 10개를 동시에 검정하므로 Bonferroni
보정 구간($\alpha' = 0.05/10$, 99.5\,\%)을 함께 보고한다. 두 계층의 상호작용은
\begin{align}
\mathrm{Int} =\;& (\texttt{llm\_unet} - \texttt{llm\_heur}) \nonumber\\
                &- (\texttt{heur\_unet} - \texttt{heur\_heur})
\label{eq:interaction}
\end{align}
로 판정한다($>0$ 상승, $\approx 0$ 가산, $<0$ 상쇄).

포획률은 체계 전체의 지표라 지휘관 계층의 기여를 가리므로, 계획 산출 시점에 확인 가능한
양을 따로 계측한다: 응답 지연, 폴백률(언어모델 실패로 휴리스틱이 대체한 비율), 커버리지
(활성 무리 중 배정된 비율), churn(재계획 간 담당을 바꾼 방어정 비율), 교차 배정(방어정--담당
무리 선분이 서로 교차하는 쌍의 수). 세 지표는 에피소드 안의 재계획에 걸쳐 평균한다.
